% Slide block for P3 (fBm and rough volatility). Frames only; included by slides.tex.
\section{Rough volatility}

\begin{frame}{fBm stopping: memory pays when increments are dependent}
\centering\tiny
\resizebox{.9\linewidth}{!}{% generated by experiments/report_tables.py from results/fbm_medium_summary.csv
\begin{tabular}{rccccc}
\toprule
$h$ & DOS & FQI-MLP & FQI-GRU & OSPG-MLP & OSPG-GRU \\
\midrule
0.05 & 0.700 (0.005) & 0.910 (0.017) & 1.297 (0.001) & 1.153 (0.004) & 1.286 (0.004) \\
0.10 & 0.508 (0.024) & 0.740 (0.004) & 1.040 (0.004) & 0.926 (0.003) & 1.032 (0.008) \\
0.20 & 0.304 (0.002) & 0.481 (0.009) & 0.643 (0.009) & 0.578 (0.004) & 0.639 (0.002) \\
0.30 & 0.187 (0.007) & 0.299 (0.008) & 0.366 (0.005) & 0.334 (0.001) & 0.362 (0.008) \\
0.40 & 0.092 (0.021) & 0.132 (0.008) & 0.158 (0.012) & 0.141 (0.010) & 0.151 (0.015) \\
0.50 & -0.001 (0.007) & -0.002 (0.009) & -0.006 (0.003) & -0.004 (0.006) & -0.008 (0.011) \\
0.60 & 0.070 (0.003) & 0.081 (0.013) & 0.099 (0.010) & 0.096 (0.003) & 0.094 (0.007) \\
0.70 & 0.165 (0.002) & 0.177 (0.000) & 0.202 (0.001) & 0.191 (0.002) & 0.193 (0.002) \\
0.80 & 0.245 (0.005) & 0.258 (0.004) & 0.260 (0.001) & 0.266 (0.004) & 0.262 (0.002) \\
0.90 & 0.302 (0.007) & 0.306 (0.003) & 0.282 (0.013) & 0.314 (0.008) & 0.306 (0.007) \\
0.95 & 0.353 (0.001) & 0.346 (0.004) & 0.328 (0.028) & 0.357 (0.002) & 0.349 (0.001) \\
\bottomrule
\end{tabular}
}
\vspace{4pt}
\begin{itemize}\scriptsize
\item Reward $R_j=W^h_{t_j}$; exact fBm by Cholesky (terminal variance $1.0022$ vs $1$ at $h=0.3$).
  20,000 paths, 100 dates, two splits, shortened training; mean (sd across splits).
\item $h\le0.3$: GRU policies clearly ahead of MLP; $h\ge0.4$: mixed. At $h=0.5$ all methods
  are close to zero, as optional stopping predicts.
\end{itemize}
\end{frame}

\begin{frame}{Rough Bergomi: model and hybrid simulation}
\small
\begin{align*}
 dS_t &= rS_t\,dt+\sqrt{v_t}S_t\,dZ_t,\qquad
 dZ_t=\rho\,dW_t+\sqrt{1-\rho^2}\,dW_t^\perp,\\
 v_t &=\xi_0\exp\!\left(\eta\widetilde W_t^h-\tfrac12\eta^2t^{2h}\right),\\
 \widetilde W_t^h&=\sqrt{2h}\int_0^t(t-s)^{h-1/2}\,dW_s.
\end{align*}
\begin{itemize}
\item Hybrid scheme: conditionally sample the near-diagonal kernel integral and approximate
  the remaining Volterra terms with hybrid evaluation weights.
\item Log-Euler stock update uses the left-point variance; put reward is
  $e^{-rt_j}(K-S_{t_j})^+$.
\item Preliminary European call check: $0.07886\pm0.00061$ (MC standard error),
  against reference $0.0791$; difference is about 0.39 standard errors.
\end{itemize}
\end{frame}

\begin{frame}{E4: American put under rough Bergomi}
\centering\scriptsize
% generated by experiments/report_tables.py from results/e4_h_sweep_summary.csv
\begin{tabular}{rcccc}
\toprule
$h$ & LSM & DOS & OSPG-MLP & OSPG-GRU \\
\midrule
0.07 & 7.935 (0.340) & 7.985 (0.350) & 8.136 (0.340) & 7.932 (0.356) \\
0.10 & 7.698 (0.329) & 7.776 (0.348) & 7.934 (0.340) & 7.690 (0.359) \\
0.30 & 7.362 (0.337) & 7.435 (0.350) & 7.507 (0.332) & 7.367 (0.368) \\
0.50 & 7.540 (0.314) & 7.619 (0.315) & 7.605 (0.352) & 7.479 (0.360) \\
\bottomrule
\end{tabular}

\vspace{5pt}
\begin{itemize}\small
\item 3,000 paths, 25 exercise intervals and one random split per $h$; shortened training.
\item OSPG-MLP has the largest point estimate in these pilot runs, but the standard errors
  overlap substantially; no method is shown to be statistically superior, and there is no
  GRU advantage.
\item Larger multi-split experiments are needed for stronger performance claims.
\end{itemize}
\end{frame}
